\begin{definition}[$\Theta_l(E)$]
  Let $E$ be a metric space equipped with its Borel $\sigma$-algebra and $l \in \mathbb{N}$. Define
  \[
    \Theta_l(E) := \{\, \gamma \in \Theta(E) : \mathbb{M}([[\gamma]]) = l \,\}
    ,
  \]
  the set of curves (\noderef{Curve}) whose current $[[\gamma]]$ (\noderef{curveCurrent}) has mass
  (\noderef{massOfFunctional}) exactly $l$. This is exactly paper.tex line 1178.

  \paragraph{Note on $\Theta_l(E)$ vs.\ length.} Membership in $\Theta_l(E)$ is a condition on the
  \emph{mass} $\mathbb{M}([[\gamma]])$ of $\gamma$'s current, not on $\length(\gamma)$; the paper
  notes (paper.tex line 1155) that $\mu_{[[\gamma]]}$ and $\mu_\gamma$, and hence
  $\mathbb{M}([[\gamma]])$ and $\length(\gamma)$, may differ in general (e.g.\ for a curve that
  traverses the same segment twice in opposite directions). Consequently $\gamma \in \Theta_l(E)$
  does \emph{not} by itself give $\length(\gamma) = l$; where this is needed (as in
  \noderef{PSFamily}), it must be recorded as a separate hypothesis.
\end{definition}
